\documentclass[10pt]{article}
 
% basic packages
\usepackage[margin=1in]{geometry}
\usepackage[pdftex]{graphicx}
\usepackage{amsmath,amssymb,amsthm}
\usepackage{custom}
\usepackage{comment}
\usepackage{enumitem}
\usepackage{multicol}
\usepackage{booktabs}
\usepackage{tabularx}

% page formatting
\usepackage{fancyhdr}
\pagestyle{fancy} 

\begin{document}
\title{Reading List for 2026}
\date{Last Updated On: Feb 27th 2026}
\maketitle

\fancyhf{}
\fancyhead[L]{2026 Reading List}
\fancyhead[R]{\leftmark}
\fancyhead[C]{\thepage}

% Making sure the right header is NOT in all caps %
\renewcommand{\sectionmark}[1]{\markright{\textsf{\arabic{section}: #1}}}
\rhead{Section\ \nouppercase{\rightmark}}
\chead{}
\cfoot{}

\begin{comment}
   ------ THESE ARE THE PREVIOUS COMMANDS FOR THE HEADER, FOOTER, ETC -------
\renewcommand{\sectionmark}[1]{\markright{\textsf{\arabic{section}. #1}}}
\renewcommand{\subsectionmark}[1]{}
\lhead{\scriptsize\textbf{\thepage}\ \ \nouppercase{\rightmark}}
\chead{}
\rhead{\normalfont MAP2302 - Spring 2026}
\lfoot{}
\cfoot{}
\rfoot{}
\setlength{\headheight}{5pt}
   ------ THESE ARE THE PREVIOUS COMMANDS FOR THE HEADER, FOOTER, ETC -------
\end{comment}

\linespread{1.03} % give a little extra room
\setlength{\parindent}{0pt}
\setlength{\parskip}{6pt}
\setcounter{secnumdepth}{2} % Numbered subsections  
\setcounter{tocdepth}{2} % numbered subsubsections in ToC

\tableofcontents

%%%%% COMPUTER SCIENCE LITERATURE %%%%%
\section{Computer Science}
This section encompasses my reading in Computer Science, 
spanning both technical and narrative works. Topics include, but are not limited to:
\begin{enumerate}
    \item Collected Works 
    \item Fiction and biographical accounts of notable and up and coming computer scientists
    \item Conceptual textbooks ranging from foundational to advanced topics, such as: \\
\end{enumerate}

\underline {Data Structures \& Algorithms }

\textbf{Data Structures can be defined as the medium through which data and is stored to be 
accessed later on.} Simply put, a data structure is a collection of data values, the relationships that take place among them, and 
the functions or operations that can be applied to the data. 

Examples of common data structures include, but are not limited to:

\begin{multicols}{2}
    \begin{itemize}[label=\textemdash, noitemsep]

        \item [] \textbf{Primitive Types:}
        \begin{enumerate}[label=\arabic*., noitemsep]
            \item Characters  
            \item Floating-Point
            \item Integer
        \end{enumerate}

        \vspace{4pt}
        \item [] \textbf{Composite Types or Non-Primitive:}
        \begin{enumerate}[label=\arabic*., noitemsep]
            \item Arrays
            \item Record ("struct")
            \item Strings 
            \item Unions 
            \item Tagged Union
        \end{enumerate}

        \columnbreak
        
        \item [] \textbf{Abstract Data Types:}
        \begin{enumerate}[label=\arabic*., noitemsep]
            \item Container 
            \item List
            \item Tuples
            \item Associative Array, Map 
            \item Set
            \item Multiset (Bag) 
            \item Stack 
            \item Queue
            \item Graph (Example: Tree, Heap)
        \end{enumerate}

    \end{itemize}
\end{multicols}  

\textbf{Algorithms} are present in the field of Mathematics and Computer Science, and algorithm is a finite (must terminate eventually) sequence 
of instructions used to solve a variety (or class) of specific problems or to perform a computation. \\

In general, a program is classified as an algorithm only if it stops eventually (satisfies the finite condition), although infinite loops, may also be considered and looked into. Mathematical logician 
George Boolos and probability theorists Richard Jeffrey define an algorithm as an explicit set of instructions for determining an output, that can be followed by a computing machine 
or human who could only carry out specific elementary operations on symbols. 
% REMAINING TOPICS TO BE EDITED %

\underline{Additional Topics Covered}
\begin{itemize}
    \item \textbf{Computer Organization and Design } \\
    \item \textbf{Operating Systems}
    \item \textbf{Discrete Structures}
    \item \textbf{Machine Learning \& Neural Networks }
    \item \textbf{Reinforcement Learning }
    \item \textbf{Natural Language processing (NLP)}
    \item \textbf{Computer Vision }
    \item \textbf{Information Theory}
\end{itemize}

\newpage
\subsection{Nonfiction}
\begin{center}
\begin{tabularx}{\textwidth}{X l l c c}
    \toprule
    \textbf{Title} & \textbf{Author} & \textbf{Level} & \textbf{Status} \\
    \midrule
    Code                                                & Charles Petzold            & Foundational               & $\square$ \\
    The Innovators                                      & Walter Issacson            & Foundational               & $\square$ \\
    How Google Works                                    & Eric Schmidt               & Foundational               & $\square$ \\
    The Pragmatic Programmer                            & A. Hunt \& D. Thomas       & Intermediate               & $\square$ \\
    Clean Code                                          & Robert C. Martin           & Intermediate               & $\square$ \\
    Machanical Intelligence                             & A.M. Turing                & Intermediate/ Advanced     & $\square$ \\
    The Art Of Doing Science and Engineering            & R. Hamming                 & Advanced                   & $\square$ \\                
    \bottomrule
\end{tabularx}
\end{center}


\subsection{Research Literature}

\subsection{Textbooks}
% Utilized tabularx instead of tabular for a flexible autofill %
\begin{center}
\begin{tabularx}{\textwidth}{X l l c c}
    \toprule
    \textbf{Title} & \textbf{Author} & \textbf{Level} & \textbf{Status} \\
    \midrule
    Neural Networks from Scratch in Python              & H. Kinsley                 & Intermediate               & $\square$ \\
    Python Data Science Handbook                        & Jake VanderPlas            & Intermediate               & $\square$ \\
    Design for How People Think                         & John Whalen                & Intermediate               & $\square$ \\
    Hands-On Machine Learning w/ Scikit-Learn           & A. Géron                   & Intermediate               & $\square$ \\
    An Introduction to Statistical Learning             & James et al.               & Intermediate / Advanced    & $\square$ \\
    Mechanical Intelligence                             & A.M. Turing                & Intermediate / Advanced    & $\square$ \\
    Reinforcement Learning                              & Sutton \& Barto            & Advanced                   & $\square$ \\
    Deep Learning                                       & Goodfellow et al.          & Advanced                   & $\square$ \\
    Deep Learning                                       & C.M. Bishop                & Advanced                   & $\square$ \\
    \bottomrule
\end{tabularx}
\end{center}

\subsection{Fiction}
% Commented out until I find books to go in here
% Utilized tabularx instead of tabular for a flexible autofill %
\begin{comment} 
\begin{center}
\begin{tabularx}{\textwidth}{X l l c c}
    \toprule
    \textbf{Title} & \textbf{Author} & \textbf{Level} & \textbf{Status} \\
    \midrule
    The Hundred-Page ML Book  & Burkov            & Foundational & $\square$ \\
    Pattern Recognition \& ML & Bishop            & Intermediate & $\square$ \\
    Deep Learning             & Goodfellow et al. & Advanced     & $\square$ \\
    Quantum ML                & Wittek            & Graduate     & $\square$ \\
    \bottomrule
\end{tabularx}
\end{center}
\end{comment}

%%%%% MATHEMATICS LITERATURE %%%%%
\section{Mathematics}
\subsection{Nonfiction}
\subsection{Research Literature}
\subsection{Textbooks}
\subsection{Fiction}


%%%%% PHYSICS LITERATURE %%%%%
\section{Physics}
\subsection{Nonfiction}
\subsection{Research Literature}
\subsection{Textbooks}
\subsection{Fiction}

%%%%% LEISURE LITERATURE %%%%%
\section{Leisure}
\subsection{Nonfiction}
\subsection{Research Literature}
\subsection{Textbooks}
\subsection{Fiction}



\section{Books Completed w/ Time Of Completion}
\subsection{Computer Science}
\subsection{Mathematics}
\subsection{Physics}
\subsection{Leisure}
\newpage

\end{document}